% Generated draft: requires mathematical and source review.
\documentclass[11pt]{article}
\usepackage[margin=25mm]{geometry}
\usepackage{fontspec}
\setmainfont{Latin Modern Roman}
\usepackage{amsmath,amssymb,mathtools,unicode-math}
\setmathfont{Latin Modern Math}
\usepackage{enumitem,needspace,xurl,hyperref}
\setlength{\parindent}{0pt}
\setlength{\parskip}{5pt}
\setlength{\emergencystretch}{3em}
\newcommand{\N}{\mathbb{N}}
\newcommand{\Z}{\mathbb{Z}}
\newcommand{\Q}{\mathbb{Q}}
\newcommand{\R}{\mathbb{R}}
\newcommand{\C}{\mathbb{C}}
\newcommand{\F}{\mathbb{F}}
\newcommand{\A}{\mathbb{A}}
\newcommand{\PP}{\mathbb{P}}
\newcommand{\E}{\mathbb{E}}
\newcommand{\eps}{\varepsilon}
\newcommand{\dd}{\,\mathrm{d}}
\begin{document}
\section*{UnsolvedMath ID: 20002474}
\textbf{Source problem code:} catalogue entry
\subsection{Definitions and mathematical statement}
This single corpus job contains two consecutive questions, identified by the aid as Problems 11 and 12, although the extraction begins ``1. (Albert Fathi).'' They must remain distinct mathematical targets within this record. On $\mathbb T^n\times\R^n$, the first asks whether a superlinear Hamiltonian, without momentum convexity, admits (i) a PDE Aubry set, for example a uniqueness set for viscosity solutions, (ii) a canonical invariant graph from the Hamiltonian dynamics, and (iii) a relation between these objects when they exist. Continuity/coercivity suffice for some PDE formulations, whereas an actual Hamiltonian flow requires additional differentiability; the source does not specify that regularity or the intended canonical constructions.

The second question defines
\[
c_0=\inf\{c:\ H(x,Du)=c\text{ admits a global viscosity subsolution on }\mathbb T^n\},
\quad
\mathcal{SS}=\{u:\ u\text{ is such a subsolution at }c_0\},
\]
and prints $S(x,y)=\sup_{u\in\mathcal{SS}}(u(x)-u(y))$. It then fixes the wrong variable for a subsolution assertion and writes the exceptional point as $0$ rather than the base point. In the corrected convention, with fixed base $a$, use
\[
S_a(x)=\sup_{u\in\mathcal{SS}}(u(x)-u(a)).
\]
With the usual continuity and uniform superlinearity hypotheses and nonempty critical subsolution class, this is the maximal critical subsolution normalized at $a$ and is a solution away from $a$. The corrected question asks whether
\[
\mathcal A_{\max}=\{a\in\mathbb T^n:S_a\text{ is a viscosity solution on all of }\mathbb T^n\}
\]
must be nonempty without convexity. The literal printed $S(x,\cdot)$ assertion is not silently retained as an equivalent formula: reversing the variable also changes the sign of the tested gradient in general.

The aid reports that Gomes--Mitake--Tran gave a one-dimensional coercive counterexample to this corrected nonemptiness question in 2018. It supplies a superlinear extension: retain the example's kinetic part $G(r)$ for $0\leq r\leq3$ and replace its outer linear tail by
\[
\widetilde G(r)=2+(r-3)+(r-3)^2\qquad(r>3).
\]
It asserts that the critical value, critical subsolutions, maximal subsolutions, and failed vertex supersolution tests are unchanged, so $\mathcal A_{\max}=\varnothing$ still holds for a superlinear nonconvex Hamiltonian. The original $G$ on $[0,3]$, spatial coefficient, regularity, and explicit failed test functions are not included, so this summary does not provide a standalone fully specified counterexample. Failure of this particular maximal-subsolution definition does not refute every possible PDE or dynamical Aubry-set construction in the first question.

A separate positive class is obtained by rescaling a Tonelli Hamiltonian about its critical level. If $H_0$ has critical value $c_0$ and
\[
\widetilde H(x,p)=c_0+b(x,p)(H_0(x,p)-c_0),\qquad0<m\leq b\leq M<\infty,
\]
with sufficient smoothness, the aid asserts that the critical value and all critical viscosity subsolutions and solutions are preserved, and hence so is their projected PDE Aubry set. On the critical level,
\[
X_{\widetilde H}=b\,X_{H_0}.
\]
Thus an invariant lifted Aubry graph of $H_0$ is preserved as a set with positive time reparametrization. Such rescalings can be genuinely nonconvex. This construction supplies a special class with inherited PDE/dynamical structure, not a general equivalence theorem for arbitrary nonconvex Hamiltonians.
\subsection{Short English statement}
The merged record asks both for nonconvex Aubry-set notions and for nonemptiness of one maximal-subsolution vertex set. The aid reports a counterexample to the corrected second question, with a superlinear tail modification, and a positive class obtained by rescaling Tonelli Hamiltonians at the critical level. Neither result settles all possible notions in the broader first question.
\subsection{Sources}
\begin{itemize}
\item Local UnsolvedMath registry, permanent ID $20002474$, \texttt{catalogue entry}; \texttt{aim-pdes-notes.json}, \texttt{notes[85]}, extracted source item 1.
\item American Institute of Mathematics, \emph{New connections between dynamical systems and PDE's}, 2003 list, Albert Fathi contribution and consecutive maximal-subsolution question, identified by the aid as Problems 11 and 12: \url{https://aimath.org/WWN/dynpde/dynpde.pdf}. Isolated $5$ and $1$ and merged ``12'' are extraction/page artifacts, distinct from the substantive variable errors.
\item Gomes--Mitake--Tran, reported 2018 counterexample to the corrected maximal-subsolution question. No title, theorem number, URL, or full authoritative bibliographic citation is supplied.
\item Research attempt 1: partial result, problem partially solved at run, rigor label ``proved,'' novelty confidence low. Registry classification partial progress, status partially solved, statement status exact. The verification field identifies a merger of two questions, without a full corrected transcription.
\end{itemize}
\end{document}
